\begin{table}[t]
\centering
\caption{The three development archetypes ($k=3$, bootstrapped from $B=200$ resamples).}
\label{tab:archetypes}
\small
\begin{tabular}{lrrrrrrrll}
\toprule
\textbf{archetype} & \textbf{n} & \textbf{mean_index} & \textbf{high} & \textbf{upper_mid} & \textbf{lower_mid} & \textbf{low} & \textbf{gaps} & \textbf{strengths} & \textbf{examples} \\
\midrule
Advanced & 66 & 4.006 & 61 & 5 & 0 & 0 & Unemployment, Forest cover, Water withdrawal & Fixed broadband, Internet users & Malaysia, Russian Federation, Greece, Switzerland, Netherlands \\
Mature middle & 73 & 0.196 & 7 & 52 & 12 & 2 & Fixed broadband, CO2 emissions per capita, Internet users & Renewable energy share, Electricity access & Kiribati, Guyana, Ecuador, Grenada, New Caledonia \\
Emerging frontier & 59 & -4.724 & 0 & 1 & 35 & 23 & Electricity access, Secondary enrolment, Life expectancy & Unemployment, Forest cover & Central African Republic, Burkina Faso, Angola, Senegal, Tajikistan \\
\bottomrule
\end{tabular}
\vspace{2pt}\parbox{0.92\linewidth}{\footnotesize \texttt{gaps}: the three indicators on which the archetype sits furthest below the global mean after orienting every indicator by the sign of its index loading, so that higher means more developed; \texttt{strengths}: the two on which it sits furthest above. The orientation is statistical, not normative: CO$_2$ emissions per capita and renewable energy share load with the income gradient, so a 'gap' there is not a policy deficit. \texttt{examples}: members sampled across the index range.}
\end{table}
